\documentclass[dvipdfmx,disablejfam]{jarticle}
\usepackage[dvipdfmx]{graphicx}
\usepackage{url}
\usepackage{amsmath}
\usepackage{amssymb}
\usepackage{bm}
\usepackage{float}
\usepackage{xcolor}

% 体裁（必要なら調整）
\topmargin -0.4mm
\oddsidemargin -0.4mm
\evensidemargin -0.4mm
\headheight 0mm
\headsep 0mm
\footskip 10mm
\textheight 247mm
\textwidth 170mm
\setlength{\parindent}{1zw}
\pagestyle{empty}

\makeatletter
\def \fgcaption{\def \@captype{figure} \caption}
\def \tbcaption{\def \@captype{table} \caption}
\makeatother

\def\figurename{Fig.}
\def\tablename{Table}
\def\topfraction{.95}
\def\bottomfraction{.95}
\def\textfraction{.05}
\def\floatpagefraction{.8}
\def\dbltopfraction{.95}
\def\dbltextfraction{.05}
\def\dblfloatpagefraction{.8}

\begin{document}

%==============================
%  ヘッダー
%==============================
\begin{center}
\begin{tabular}{|c|c|c|} \hline
  \parbox{50mm}{\vspace{1mm} \centering{\Large\bf 2026/04/30}} &
  \parbox{55mm}{\vspace{1mm} \centering{\Large\bf M1\quad 市川 詠祥}} &
  \parbox{50mm}{\vspace{1mm} \centering{\Large\bf No.~04}} \\ [2mm] \hline

  \multicolumn{3}{|c|}{
    \parbox{155mm}{
      \centering\Large\bf \vspace{2mm}
       円形きずの応力分布プログラム作成について
    }
  } \\ [3mm] \hline\hline

  \multicolumn{3}{|c|}{
    \parbox{155mm}{
      \large\vspace{2mm}

      {\Large\bf 今月の目標}
      \begin{itemize}
        \item 円形きず（U字型）の応力可視化プログラムの作成
        \item 各計測点でのT1・T3応力分布の確認
      \end{itemize}

      {\Large\bf 進捗の要点}
      \begin{itemize}
        \item 円形きず（U字型）の応力可視化プログラム（202604\_ushape）を作成した
        \item ピッチ$P=2.00$\,mm・深さ$d=0.20$\,mmの条件でシミュレーションを実施し，T1・T3応力分布を可視化した
        \item 各コーナー計測点（Pt\,1〜Pt\,9）におけるT1応力の時間遷移を比較した結果，
              Point\,1とPoint\,9においてT1応力が他の計測点に比べて著しく小さくなった
        \item 原因について調査中
      \end{itemize}

      {\Large\bf 今後の実施事項}
      \begin{itemize}
        \item Point\,1・Point\,9のT1応力が小さい原因の解明
        \item 
      \end{itemize}

      {\Large\bf 投稿・発表予定}
      \begin{itemize}
        \item 5/15 , 5/22, 6/12, 6/19 のうち2回, 輪講での発表（PINN関連の内容を予定）
      \end{itemize}
    }
  } \\ [3mm] \hline
\end{tabular}
\end{center}

\vspace{6mm}
{\Large\bf 今回の内容}
\vspace{3mm}

\subsection*{1. 円形きず解析プログラムの作成}

前回までに作成したくさび型きずのプログラムをベースとして，
円形きず（U字型）を対象とした解析プログラム群（202604\_ushape）を新規作成した．
U字型きずの形状は，直線部（深さ$d - R$）と半円（半径$R = w/2$）からなり，
ピッチはきず + 平らな部分で定義される．

計測点はFig.\,\ref{fig:map_points}に示すように，溝右側の弧上のコーナー点をPt\,1〜Pt\,8，
底面に一番近いところ（$x = 1999$, 底面境界の1格子上）をPt\,9として設定した．

\begin{figure}[H]
  \centering
  \includegraphics[width=0.85\textwidth]{../../project4/tmp_output/ushape_map_points_pitch200_width25_depth20.png}
  \caption{U字型きずの計測点配置（$P=2.00$\,mm, $w=0.25$\,mm, $d=0.20$\,mm）}
  \label{fig:map_points}
\end{figure}

\subsection*{2. T1・T3応力分布の確認}

$P=2.00$\,mm，$d=0.20$\,mmの条件でシミュレーションを実施した結果を
Fig.\,\ref{fig:surface_map}に示す．
上段はコーナー計測点Pt\,1〜Pt\,9における時間方向のT1応力，
下段はすきま部分(z = 20.25mm ～ 22.00mm)のT3応力の分布である．

\begin{figure}[H]
  \centering
  \includegraphics[width=0.85\textwidth]{../../project4/tmp_output/ushape_T1_T3_surface_map_center_pitch200_depth20.png}
  \caption{T1・T3応力の分布（$P=2.00$\,mm, $d=0.20$\,mm）}
  \label{fig:surface_map}
\end{figure}

\subsection*{3. Point\,1・Point\,9のT1応力が小さい問題}

各計測点のT1応力の時間波形をFig.\,\ref{fig:waveform}に示す．
Point\,1と, Point\,9（底面境界のすぐ上）において，他の計測点と比較してT1応力が著しく小さいことが確認された．

\begin{figure}[H]
  \centering
  \includegraphics[width=0.85\textwidth]{../../project4/tmp_output/T1_waveform_all_points_pitch200_depth20.png}
  \caption{各計測点におけるT1応力の時間波形（$P=2.00$\,mm, $d=0.20$\,mm）}
  \label{fig:waveform}
\end{figure}

Point\,1とPoint\,9のT1応力が小さい理由について, 原因を調べています. 

\subsection*{4. 今後の予定}

\begin{itemize}
  \item Point\,1とPoint\,9のT1応力が小さい原因の特定
  \item 輪講にむけた資料の作成（PINN関連の内容にしようと思います.）
\end{itemize}

%===========================================

\end{document}
